\chapter{Guía de usuario}

\section{Acceso y sesión}

Abra \href{\FrontendUrl}{\FrontendUrl}. Los clientes nuevos completan el registro y confirman su
correo antes de iniciar sesión. Los trabajadores y administradores reciben sus cuentas por el canal
definido por el responsable operativo.

Si una sesión expira, la aplicación intenta renovarla mediante una cookie segura. Si no puede,
solicita un nuevo inicio de sesión. No comparta contraseñas ni utilice cuentas demo en producción.

\section{Flujo del Cliente}

\begin{enumerate}
  \item Inicie sesión y abra el panel de tickets.
  \item Seleccione \textbf{Crear ticket}.
  \item Elija el servicio y la prioridad; escriba asunto y descripción suficiente para reproducir el
  problema.
  \item Adjunte únicamente archivos necesarios y sin secretos.
  \item Envíe y conserve el número generado.
  \item Consulte el detalle para revisar estado, responsable, comentarios e historial.
\end{enumerate}

El Cliente no asigna trabajadores ni modifica directamente los estados operativos. Si necesita
corregir información, debe añadir un comentario o contactar al responsable indicado.

\section{Flujo del Trabajador}

\begin{enumerate}
  \item Revise los tickets asignados y filtre por estado o prioridad.
  \item Abra el detalle antes de actuar; confirme servicio, evidencia e historial.
  \item Registre avances como comentarios concretos y sin datos sensibles innecesarios.
  \item Cambie el estado e incluya el motivo. Un comentario vacío es rechazado.
  \item Marque como resuelto o cerrado solo cuando exista evidencia operativa suficiente.
\end{enumerate}

La reapertura está permitida. Debe explicar la razón para mantener una línea de tiempo útil.

\section{Flujo del Administrador}

El Administrador puede ver el conjunto de tickets, asignar o reasignar, gestionar usuarios,
catálogo, galería, logos y testimonios, y consultar reportes.

Para asignar un ticket:

\begin{enumerate}
  \item abra un ticket \textbf{Nuevo};
  \item seleccione un trabajador activo;
  \item registre el comentario de asignación;
  \item confirme que el estado cambió a \textbf{EnProceso} y que el evento aparece en el historial.
\end{enumerate}

Para bloquear un usuario, registre primero la razón operativa y confirme el impacto sobre sesiones
activas. No elimine evidencia de un incidente en curso.

\section{Notificaciones y reportes}

El panel de notificaciones permite consultar y marcar eventos como leídos y ajustar preferencias.
La ausencia temporal de WebSocket no impide usar la API; la actualización puede requerir recargar.

Los reportes administrativos admiten filtros y exportación en \textbf{PDF} o \textbf{Excel}. Antes
de compartir un archivo, revise nombres, correos, comentarios y cualquier otro dato personal.

\section{Problemas frecuentes}

\begin{tabularx}{\textwidth}{@{}L{0.28\textwidth}Y@{}}
\toprule
\textbf{Síntoma} & \textbf{Acción recomendada}\\
\midrule
No inicia sesión & Verifique correo confirmado, credenciales y estado de la cuenta; luego revise la API.\\
El backend tarda & Espere el despertar del servicio y compruebe \href{\ApiHealthUrl}{el health check}.\\
No hay tiempo real & Recargue, valide `VITE\_WS\_URL`, origen permitido y disponibilidad de Redis.\\
No llega correo & Revise spam, remitente, proveedor y logs; no repita el registro de forma masiva.\\
Adjunto rechazado & Compruebe tamaño y tipo; no desactive la validación para forzar una carga.\\
Exportación vacía & Revise filtros, periodo y permisos del usuario administrador.\\
\bottomrule
\end{tabularx}
